\documentclass[11pt]{article}
\usepackage{amsmath,amssymb,amsthm}
\usepackage[margin=1.1in]{geometry}
\newtheorem{theorem}{Theorem}
\newtheorem{lemma}[theorem]{Lemma}
\newcommand{\vtwo}{v_2}

\title{Erd\H{o}s Problem \#404: the $p=2$ classification\\ \large The obstruction families $a\equiv 4 \pmod 8$, $a\equiv 10 \pmod{16}$, $a\equiv 32,48,8 \pmod{64}$, $a\equiv 40,56,88,70,118 \pmod{128}$}
\date{Campaign Kills \#11--\#13 --- kernel-green, CI-verified}

\begin{document}
\maketitle

\section*{The unsolved problem (Erd\H{o}s--Graham 1980, p.\,79; erdosproblems.com/404)}

\textbf{Question 1.} For which integers $a \ge 1$ and primes $p$ is there a
finite upper bound on those $k$ such that there are
$a = a_1 < a_2 < \cdots < a_n$ with $p^k \mid (a_1! + a_2! + \cdots + a_n!)$?

\textbf{Question 2.} If $f(a,p)$ is the greatest such $k$, how does $f$ behave?

\textbf{Question 3.} Is there a prime $p$ and an infinite sequence
$a_1 < a_2 < \cdots$ such that $v_p\!\left(\sum_{i \le k} a_i!\right) \to \infty$?

The problem page records (Li76): $f(2,2) \le 2^{54}$.

\section*{The reduction (campaign, kernel-green in \texttt{PadicReduction.lean})}

Write $P_j(a) = (a+1)(a+2)\cdots(a+j) = \dfrac{(a+j)!}{a!}$ for $j \ge 1$.
Every admissible sum factors as
\[
a! + \sum_{i \in B} i! \;=\; a! \left( 1 + \sum_{j \in J} P_j(a) \right),
\qquad J = \{i - a : i \in B\},
\]
so by additivity of $\vtwo$,
\begin{equation}
f(a,2) = \vtwo(a!) + g(a,2), \qquad
g(a,2) = \max\Bigl\{ g : \exists\, J \subseteq \{1,2,3,\dots\},\;
2^g \,\Big|\, 1 + \textstyle\sum_{j \in J} P_j(a) \Bigr\}.
\label{eq:reduction}
\end{equation}

\section*{The theorems}

\begin{theorem}[T2: the $a \equiv 4 \pmod 8$ family]
For every $a \equiv 4 \pmod 8$:
\begin{enumerate}
\item[(i)] \emph{(Witness.)} $1 + P_1(a) + P_3(a) = (a+1) + (a+1)(a+2)(a+3) + 1
\equiv 0 \pmod 8$ but $\not\equiv 0 \pmod{16}$; hence $g(a,2) \ge 3$.
\item[(ii)] \emph{(Obstruction.)} For every $J \subseteq \{1,2,\dots\}$:
$1 + \sum_{j \in J} P_j(a) \not\equiv 0 \pmod{16}$; hence $g(a,2) \le 3$.
\end{enumerate}
Therefore $f(a,2) = \vtwo(a!) + 3$ exactly, on the whole family.
\end{theorem}

\begin{proof}[Proof of (ii)]
Since $a \equiv 4 \pmod 8$: $\vtwo(a+2) = 1$ and $\vtwo(a+4) = 3$, so every
$P_j(a)$ with $j \ge 4$ contains the factor $(a+2)(a+4)$ and is
$\equiv 0 \pmod{16}$. The sum $1 + \sum_J P_j$ is therefore congruent mod 16
to a sub-sum of the three \emph{window coins}
\[
P_1(a) = a+1,\quad P_2(a) = (a+1)(a+2),\quad P_3(a) = (a+1)(a+2)(a+3).
\]
Case $a \equiv 4 \pmod{16}$: the coin residues are
$\{5, 14, 2\}$, whose eight sub-sums are
$\{0, 2, 3, 5, 7, 14\}$ --- none is $\equiv 15 \pmod{16}$.
Case $a \equiv 12 \pmod{16}$: the residues are $\{13, 6, 10\}$, sub-sums
$\{0, 3, 6, 7, 10, 13\}$ --- again $15$ is missed.
In both cases $1 + \sum \not\equiv 0 \pmod{16}$.
\end{proof}

\begin{theorem}[T3: the $a \equiv 10 \pmod{16}$ family]
For every $a \equiv 10 \pmod{16}$: the witness $J = \{1,3\}$ gives
$1 + P_1 + P_3 \equiv 0 \pmod{64}$ exactly, while every $J$ has
$1 + \sum_J P_j \not\equiv 0 \pmod{128}$. Hence
$f(a,2) = \vtwo(a!) + 6$ exactly.
\end{theorem}

\begin{proof}
Here $\vtwo(a+2)=2, \vtwo(a+4)=1, \vtwo(a+6)=4$, so $2^2 \cdot 2^1 \cdot 2^4
= 128$ divides every $P_j(a)$ with $j \ge 6$: the window is the five coins
$P_1,\dots,P_5$. Modulo 128 the window residues depend only on
$a \bmod 128 \in \{10, 26, 42, 58, 74, 90, 106, 122\}$; in each of the eight
cases the $32$ sub-sums of the five coins miss $-1$ (machine-verified and
\emph{decided inside the Lean kernel} as eight boolean leaves). The witness
$1 + (a+1) + (a+1)(a+2)(a+3) \equiv 0 \pmod{64}$ holds identically on the
family (four sub-cases mod 64, decided).
\end{proof}

\begin{theorem}[T4: the $a \equiv 32, 48 \pmod{64}$ family]
For every $a \equiv 32 \pmod{64}$ the witness $J = \{1,3,5\}$, and for every
$a \equiv 48 \pmod{64}$ the witness $J = \{1,3,4\}$, attain
$g = 6$ exactly; every $J$ misses $-1$ mod 128. Hence
$f(a,2) = \vtwo(a!) + 6$ exactly on the family.
\end{theorem}

\begin{proof}
$\vtwo(a+2) = \vtwo(a+6) = 1$, $\vtwo(a+4) = 2$, $\vtwo(a+8) = 3$, so
$2\cdot 4 \cdot 2 \cdot 8 = 128$ divides every $P_j(a)$ with $j \ge 8$:
the window is the seven coins $P_1, \dots, P_7$. Four residue classes
mod 128 ($\{32, 96\}$ and $\{48, 112\}$); in each, all $128$ sub-sums miss
$-1$ (decided in the kernel).
\end{proof}

\begin{theorem}[T5: the $a \equiv 8 \pmod{64}$ family]
For every $a \equiv 8 \pmod{64}$ the witness $J = \{1,3,4,6\}$ attains
$g = 6$ exactly; every $J$ misses $-1$ mod 128. Hence
$f(a,2) = \vtwo(a!) + 6$ exactly on the family.
\end{theorem}

\begin{proof}
$\vtwo(a+2) = \vtwo(a+6) = 1$, $\vtwo(a+4) = 2$, $8 \mid a+8$, so
$2\cdot 4 \cdot 2 \cdot 8 = 128$ divides every $P_j(a)$ with $j \ge 8$:
the window is the seven coins $P_1, \dots, P_7$. Two residue classes
mod 128 ($\{8, 72\}$): the window residues are
$\{9, 90, 94, 104, 72, 112, 16\}$ and $\{73, 26, 30, 104, 72, 112, 16\}$;
in each, all $128$ sub-sums miss $-1$ (decided in the kernel as two
boolean leaves). The witness $1 + P_1 + P_3 + P_4 + P_6 \equiv 0
\pmod{64}$ holds identically on the family (single sub-case mod 64,
decided; $v_2$ of the factor is exactly 6 since the sum is
$\equiv 64 \pmod{128}$).
\end{proof}

\section*{The five new families (Kill \#13): T6/T6b/T6c/T7/T7b}

\begin{theorem}[T6: the $a \equiv 40 \pmod{128}$ family]
For every $a \equiv 40 \pmod{128}$ the witness $J = \{1,3,4,6\}$ attains
$g = 7$ exactly; every $J$ misses $-1$ mod 256. Hence
$f(a,2) = \vtwo(a!) + 7$ exactly on the family.
\end{theorem}

\begin{theorem}[T6b: the $a \equiv 56 \pmod{128}$ family]
For every $a \equiv 56 \pmod{128}$ the witness $J = \{1,3,5,6\}$ attains
$g = 7$ exactly; every $J$ misses $-1$ mod 256. Hence
$f(a,2) = \vtwo(a!) + 7$ exactly on the family.
\end{theorem}

\begin{theorem}[T6c: the $a \equiv 88 \pmod{128}$ family]
For every $a \equiv 88 \pmod{128}$ the witness $J = \{1,3,4,7\}$ attains
$g = 7$ exactly; every $J$ misses $-1$ mod 256. Hence
$f(a,2) = \vtwo(a!) + 7$ exactly on the family.
\end{theorem}

\begin{proof}[Proof of T6/T6b/T6c]
The three families share one vanish profile:
$\vtwo(a+2) = 1$, $\vtwo(a+4) = 2$, $\vtwo(a+6) = 1$, $\vtwo(a+8) = 4$
(that is, $2 \mid a+2$, $4 \mid a+4$, $2 \mid a+6$, $16 \mid a+8$), so
$2^{1+2+1+4} = 256$ divides every $P_j(a)$ with $j \ge 8$: the window is
the seven coins $P_1, \dots, P_7$ (the lemma \texttt{vanish\_T6fam}, proved
by $2$-adic factor extraction and induction). Six residue classes mod 256
($\{40, 168\}$, $\{56, 184\}$, $\{88, 216\}$); in each, all $2^7 = 128$
sub-sums of the window coins miss $-1$ (decided in the kernel as six
boolean case trees \texttt{leaf\_T6\_40}~$\dots$~\texttt{leaf\_T6c\_216}),
giving the obstructions \texttt{obs\_T6/T6b/T6c} and the upper bounds
$f \le \vtwo(a!) + 7$. The three witnesses attain
$1 + \sum_{j \in J} P_j \equiv 0 \pmod{128}$ but $\not\equiv 0
\pmod{256}$ identically on their families (decided), so $g \ge 7$ and
$f = \vtwo(a!) + 7$ exactly (\texttt{f404\_T6\_exact} and siblings).
\end{proof}

\begin{theorem}[T7: the $a \equiv 70 \pmod{128}$ family]
For every $a \equiv 70 \pmod{128}$ the witness $J = \{1,3,6,7,9\}$ attains
$g = 10$ exactly; every $J$ misses $-1$ mod 2048. Hence
$f(a,2) = \vtwo(a!) + 10$ exactly on the family.
\end{theorem}

\begin{theorem}[T7b: the $a \equiv 118 \pmod{128}$ family]
For every $a \equiv 118 \pmod{128}$ the same witness $J = \{1,3,6,7,9\}$
attains $g = 10$ exactly; every $J$ misses $-1$ mod 2048. Hence
$f(a,2) = \vtwo(a!) + 10$ exactly on the family.
\end{theorem}

\begin{proof}[Proof of T7/T7b]
The two families share the vanish profile
$\vtwo(a+2) = 3$, $\vtwo(a+4) = 1$, $\vtwo(a+6) = 2$, $\vtwo(a+8) = 1$,
$\vtwo(a+10) = 4$ (that is, $8 \mid a+2$, $2 \mid a+4$, $4 \mid a+6$,
$2 \mid a+8$, $16 \mid a+10$), so $2^{3+1+2+1+4} = 2048$ divides every
$P_j(a)$ with $j \ge 10$ (\texttt{vanish\_T7fam}): the window is the NINE
coins $P_1, \dots, P_9$, and the leaf residues are the sixteen classes
$a \bmod 2048 \in \{70, 198, \dots, 1990\}$ (T7) and
$\{118, 246, \dots, 2038\}$ (T7b). In each of the 32 leaves, all
$2^9 = 512$ sub-sums miss $-1$ mod 2048; the witness attains
$1 + \sum_{j \in J} P_j \equiv 0 \pmod{1024}$ but $\not\equiv 0
\pmod{2048}$ identically (decided), whence
$f = \vtwo(a!) + 10$ exactly (\texttt{f404\_T7\_exact}, \texttt{f404\_T7b\_exact}).
\end{proof}

\begin{proof}[Proof of the 32 leaves: the memory-safe \texttt{mc9}/\texttt{boundedForall} architecture]
A naive tactic proof of a T7 leaf --- case-splitting the nine membership
booleans for each of the 32 residues --- spawns $32 \times 2^9$ goals and
exhausted the memory of the development sandbox. The kernel instead packs
each leaf into ONE closed term: \texttt{memMask}\,$(a, B)$ is the 9-bit
membership mask of the window (bit $j-1$ set iff $a+j \in B$; provably
$< 512$), \texttt{mc9}\,$(r, n)$ reconstructs the sub-sum
$\sum_{j=1}^{9} P_j(r)\,\mathrm{bit}_j(n)$ from the mask $n$, and
\[
\texttt{boundedForall}\;(\lambda n.\; (1 + \texttt{mc9}\; r\; n) \bmod 2048 \ne 0)\;512 = \texttt{true}
\]
is a single kernel-evaluated \texttt{decide} per leaf residue
(\texttt{check\_T7\_70}~$\dots$~\texttt{check\_T7b\_2038}: 32 checks, each
evaluating all 512 sub-sums). The bridge lemma \texttt{bf\_true} together
with \texttt{memMask\_lt} then instantiates the checked leaf at the actual
membership mask of an arbitrary $B$. The entire T7/T7b obstruction is thus
32 kernel decides with zero tactic explosion.
\end{proof}

\section*{The classification data}

The engine (exact bitset subset-sum DP) certifies:
\begin{itemize}
\item family uniformity: $g = 3$ on all $a \equiv 4 \pmod 8$ up to $204$;
$g = 6$ on all $a \equiv 10 \pmod{16}$ up to $202$; $g = 6$ on all
$a \equiv 32, 48 \pmod{64}$ up to $304$; $g = 6$ on all $a \equiv 8
\pmod{64}$ up to $456$ (T5); $g = 7$ on all $a \equiv 40, 56, 88
\pmod{128}$ and $g = 10$ on all $a \equiv 70, 118 \pmod{128}$, seven
members each up to $a < 900$ (T6/T6b/T6c/T7/T7b, new);
\item the small-$g$ table for all $a \le 200$ (e.g.\ $g(22,2) = 11$,
$g(80,2) = 15$, $g(120,2) = 9$);
\item the \textbf{deep survivors} at modulus $2^{28}$:
$a \in \{2, 6, 14, 16, 18, 24, 30, 34, 38, 62, 64, 78, 94, 98, 102, 126,
142, 146, 158, 162, 190,  194\}$ have $g \ge 28$;
\item bounded deep values $g(46,2) = 27$, $g(86,2) = 26$, $g(104,2) = 23$,
$g(110,2) = 21$, $g(144,2) = 18$, $g(174,2) = 23$, $g(400,2) = 14$;
\item odd primes: $g(1,3), g(4,3), g(7,3), g(9,3), g(10,3), g(12,3)
\ge 19$ at $3^{19}$ --- the odd-$p$ landscape is deep-survivor-rich and
admits no uniform obstruction at scan depth.
\end{itemize}

\section*{What remains open}

Questions 1 and 3 for the deep survivors (including $(2,2)$, where Lin's
$f(2,2) \le 2^{54}$ is the only finiteness evidence, and our $g(2,2) \ge 28$),
the odd-$p$ classification, and Question 3 in general. No $(a,p)$ with
$g(a,p) = \infty$ is known or refuted.

\section*{The Lean kernel}

\texttt{ObstructionFamilies.lean} (3338 lines, core Lean 4, zero sorry):
the nine family theorems are proved for \emph{all} $a$ in each family, with the
obstruction leaves decided by the kernel (boolean case trees over the window
coins for T2--T6c; one closed-term \texttt{decide} per leaf via the
\texttt{mc9}/\texttt{boundedForall} bridge for T7/T7b) and the uniform vanish
lemmas by $2$-adic factor extraction. Axioms:
$[\mathrm{propext}, \mathrm{Quot.sound}]$ (T2/T3),
$[\mathrm{propext}, \mathrm{Classical.choice}, \mathrm{Quot.sound}]$
(T4--T7b). Receipts (kernel \texttt{\#eval}): $f(4,2)=6$, $f(12,2)=13$,
$f(20,2)=21$, $f(36,2)=37$, $f(10,2)=14$, $f(26,2)=29$, $f(42,2)=45$,
$f(32,2)=37$, $f(48,2)=52$ --- recovering Theorem C ($f(4,2)=6$) of the
campaign's problem404-full kernel as the $a=4$ instance of T2; T5:
$f(8,2)=13$, $f(72,2)=76$, $f(136,2)=140$, $f(200,2)=203$; and the five new
families (Kill \#13): $f(40,2)=45$, $f(168,2)=172$, $f(296,2)=300$,
$f(424,2)=427$ (T6); $f(56,2)=60$, $f(184,2)=187$ (T6b); $f(88,2)=92$,
$f(216,2)=219$ (T6c); $f(70,2)=77$, $f(198,2)=204$, $f(326,2)=332$,
$f(454,2)=459$ (T7); $f(118,2)=123$, $f(246,2)=250$ (T7b).

\end{document}
